% Options for packages loaded elsewhere
\PassOptionsToPackage{unicode}{hyperref}
\PassOptionsToPackage{hyphens}{url}
%
\documentclass[
  ignorenonframetext,
  aspectratio=54,
]{beamer}
\usepackage{pgfpages}
\setbeamertemplate{caption}[numbered]
\setbeamertemplate{caption label separator}{: }
\setbeamercolor{caption name}{fg=normal text.fg}
\beamertemplatenavigationsymbolsempty
% Prevent slide breaks in the middle of a paragraph
\widowpenalties 1 10000
\raggedbottom
\setbeamertemplate{part page}{
  \centering
  \begin{beamercolorbox}[sep=16pt,center]{part title}
    \usebeamerfont{part title}\insertpart\par
  \end{beamercolorbox}
}
\setbeamertemplate{section page}{
  \centering
  \begin{beamercolorbox}[sep=12pt,center]{part title}
    \usebeamerfont{section title}\insertsection\par
  \end{beamercolorbox}
}
\setbeamertemplate{subsection page}{
  \centering
  \begin{beamercolorbox}[sep=8pt,center]{part title}
    \usebeamerfont{subsection title}\insertsubsection\par
  \end{beamercolorbox}
}
\AtBeginPart{
  \frame{\partpage}
}
\AtBeginSection{
  \ifbibliography
  \else
    \frame{\sectionpage}
  \fi
}
\AtBeginSubsection{
  \frame{\subsectionpage}
}
\usepackage{amsmath,amssymb}
\usepackage{iftex}
\ifPDFTeX
  \usepackage[T1]{fontenc}
  \usepackage[utf8]{inputenc}
  \usepackage{textcomp} % provide euro and other symbols
\else % if luatex or xetex
  \usepackage{unicode-math} % this also loads fontspec
  \defaultfontfeatures{Scale=MatchLowercase}
  \defaultfontfeatures[\rmfamily]{Ligatures=TeX,Scale=1}
\fi
\usepackage{lmodern}
\ifPDFTeX\else
  % xetex/luatex font selection
\fi
% Use upquote if available, for straight quotes in verbatim environments
\IfFileExists{upquote.sty}{\usepackage{upquote}}{}
\IfFileExists{microtype.sty}{% use microtype if available
  \usepackage[]{microtype}
  \UseMicrotypeSet[protrusion]{basicmath} % disable protrusion for tt fonts
}{}
\makeatletter
\@ifundefined{KOMAClassName}{% if non-KOMA class
  \IfFileExists{parskip.sty}{%
    \usepackage{parskip}
  }{% else
    \setlength{\parindent}{0pt}
    \setlength{\parskip}{6pt plus 2pt minus 1pt}}
}{% if KOMA class
  \KOMAoptions{parskip=half}}
\makeatother
\usepackage{xcolor}
\newif\ifbibliography
\usepackage{color}
\usepackage{fancyvrb}
\newcommand{\VerbBar}{|}
\newcommand{\VERB}{\Verb[commandchars=\\\{\}]}
\DefineVerbatimEnvironment{Highlighting}{Verbatim}{commandchars=\\\{\}}
% Add ',fontsize=\small' for more characters per line
\newenvironment{Shaded}{}{}
\newcommand{\AlertTok}[1]{\textcolor[rgb]{1.00,0.00,0.00}{\textbf{#1}}}
\newcommand{\AnnotationTok}[1]{\textcolor[rgb]{0.38,0.63,0.69}{\textbf{\textit{#1}}}}
\newcommand{\AttributeTok}[1]{\textcolor[rgb]{0.49,0.56,0.16}{#1}}
\newcommand{\BaseNTok}[1]{\textcolor[rgb]{0.25,0.63,0.44}{#1}}
\newcommand{\BuiltInTok}[1]{\textcolor[rgb]{0.00,0.50,0.00}{#1}}
\newcommand{\CharTok}[1]{\textcolor[rgb]{0.25,0.44,0.63}{#1}}
\newcommand{\CommentTok}[1]{\textcolor[rgb]{0.38,0.63,0.69}{\textit{#1}}}
\newcommand{\CommentVarTok}[1]{\textcolor[rgb]{0.38,0.63,0.69}{\textbf{\textit{#1}}}}
\newcommand{\ConstantTok}[1]{\textcolor[rgb]{0.53,0.00,0.00}{#1}}
\newcommand{\ControlFlowTok}[1]{\textcolor[rgb]{0.00,0.44,0.13}{\textbf{#1}}}
\newcommand{\DataTypeTok}[1]{\textcolor[rgb]{0.56,0.13,0.00}{#1}}
\newcommand{\DecValTok}[1]{\textcolor[rgb]{0.25,0.63,0.44}{#1}}
\newcommand{\DocumentationTok}[1]{\textcolor[rgb]{0.73,0.13,0.13}{\textit{#1}}}
\newcommand{\ErrorTok}[1]{\textcolor[rgb]{1.00,0.00,0.00}{\textbf{#1}}}
\newcommand{\ExtensionTok}[1]{#1}
\newcommand{\FloatTok}[1]{\textcolor[rgb]{0.25,0.63,0.44}{#1}}
\newcommand{\FunctionTok}[1]{\textcolor[rgb]{0.02,0.16,0.49}{#1}}
\newcommand{\ImportTok}[1]{\textcolor[rgb]{0.00,0.50,0.00}{\textbf{#1}}}
\newcommand{\InformationTok}[1]{\textcolor[rgb]{0.38,0.63,0.69}{\textbf{\textit{#1}}}}
\newcommand{\KeywordTok}[1]{\textcolor[rgb]{0.00,0.44,0.13}{\textbf{#1}}}
\newcommand{\NormalTok}[1]{#1}
\newcommand{\OperatorTok}[1]{\textcolor[rgb]{0.40,0.40,0.40}{#1}}
\newcommand{\OtherTok}[1]{\textcolor[rgb]{0.00,0.44,0.13}{#1}}
\newcommand{\PreprocessorTok}[1]{\textcolor[rgb]{0.74,0.48,0.00}{#1}}
\newcommand{\RegionMarkerTok}[1]{#1}
\newcommand{\SpecialCharTok}[1]{\textcolor[rgb]{0.25,0.44,0.63}{#1}}
\newcommand{\SpecialStringTok}[1]{\textcolor[rgb]{0.73,0.40,0.53}{#1}}
\newcommand{\StringTok}[1]{\textcolor[rgb]{0.25,0.44,0.63}{#1}}
\newcommand{\VariableTok}[1]{\textcolor[rgb]{0.10,0.09,0.49}{#1}}
\newcommand{\VerbatimStringTok}[1]{\textcolor[rgb]{0.25,0.44,0.63}{#1}}
\newcommand{\WarningTok}[1]{\textcolor[rgb]{0.38,0.63,0.69}{\textbf{\textit{#1}}}}
\setlength{\emergencystretch}{3em} % prevent overfull lines
\providecommand{\tightlist}{%
  \setlength{\itemsep}{0pt}\setlength{\parskip}{0pt}}
\setcounter{secnumdepth}{-\maxdimen} % remove section numbering
\usepackage{pgfpages}
\usepackage{luatexbase}
\usepackage{microtype}
\usepackage{luatextra}
\usepackage{luaotfload}
\usepackage{fontspec}
\usepackage{unicode-math}

\setmainfont[
Extension=.otf,
UprightFont={IBM Plex Serif-Text},
BoldFont={IBM Plex Serif-Bold},
ItalicFont={IBM Plex Serif-Italic},
BoldItalicFont={IBM Plex Serif-BoldItalic}
]{IBM Plex Serif}[Ligatures=TeX,Scale=MatchUppercase]
% Scales Plex to default size, e.g., to Helvetica´s size.

\setmonofont{IBM Plex Mono Text}[Scale=MatchLowercase,
Ligatures=TeX]
\setsansfont{IBM Plex Sans Text}[Scale=MatchUppercase,
Ligatures=TeX]
\setmathfont{Latin Modern Math}

\definecolor{CTred}{RGB}{229,32,32}
\definecolor{CTgrey}{RGB}{153,153,153}


% colors: white text on 90% black background
\setbeamercolor{normal text}{fg=black,bg=white}

% light blue as a highlight color
\setbeamercolor*{structure}{fg=CTred}
\setbeamercolor{section title}{fg=CTred}
\setbeamercolor{alerted text}{use=structure,fg=CTred}
\setbeamercolor*{palette primary}{use=structure,fg=structure.fg}
\setbeamercolor*{palette secondary}{use=structure,fg=structure.fg!95!black}
\setbeamercolor*{palette tertiary}{use=structure,fg=structure.fg!90!black}
\setbeamercolor*{palette quaternary}{use=structure,fg=structure.fg!95!black,bg=black!80}

\setbeamercolor*{framesubtitle}{fg=white}


% use system fonts: here, Gill Sans
\usefonttheme{professionalfonts}
\setbeamerfont{quote}{shape=\upshape}

% eliminate silly beamer navigation line at bottom of slides
\setbeamertemplate{navigation symbols}{}

% ensure text jusfication
\usepackage{ragged2e}
\justifying

% pandoc makes 2nd-lever headers into blocks, and this ensures justification
% in blocks too
\addtobeamertemplate{block begin}{}{\justifying}




\urlstyle{same}
\usepackage[overlay,absolute]{textpos}

\setbeamertemplate{items}[square]

\TPGrid[10 mm,8 mm]{9}{8}
% beamer's left and right margin is 10 mm. The top/bottom margin is ??
% or without a header ??
% the slide dimensions are 128 mm x 96 mm
% so the resulting \TPHorizModule = 12 mm and \TPVertModule = 10 mm

% uncomment if you want biblatex for citations on slides

% \usepackage{csquotes}
% \usepackage[notes,short,noibid,backend=biber]{biblatex-chicago}
% \bibliography{course.bib} 

\titlegraphic{\includegraphics[width=0.33\paperwidth]{../assets/logo/Horizontal/RGB/pdf/Koren_logo_horizontal_RGB_colorful.pdf}}
\providecommand{\exhibit}[2]{\includegraphics[keepaspectratio, height=0.9\textheight, width=\textwidth]{assets/img/#1}\\ {\tiny #2}}
\ifLuaTeX
  \usepackage{selnolig}  % disable illegal ligatures
\fi
\IfFileExists{bookmark.sty}{\usepackage{bookmark}}{\usepackage{hyperref}}
\IfFileExists{xurl.sty}{\usepackage{xurl}}{} % add URL line breaks if available
\urlstyle{same}
\hypersetup{
  pdftitle={Command Line and Version Control},
  pdfauthor={Miklós Koren (koren.mk)},
  hidelinks,
  pdfcreator={LaTeX via pandoc}}

\title{Command Line and Version Control}
\author{Miklós Koren (koren.mk)}
\date{}

\begin{document}
\frame{\titlepage}

\begin{frame}{Tools}
\protect\hypertarget{tools}{}
\begin{block}{Git Bash}
\protect\hypertarget{git-bash}{}
https://gitforwindows.org/
\end{block}

\begin{block}{Dataset}
\protect\hypertarget{dataset}{}
Second at https://thnk.ng/day1
\end{block}

\begin{block}{Video recording}
\protect\hypertarget{video-recording}{}
https://thnk.ng/2023-10-12
\end{block}
\end{frame}

\begin{frame}{Lessons from the Command Line}
\protect\hypertarget{lessons-from-the-command-line}{}
\begin{enumerate}
\tightlist
\item
  Everything is a file.
\item
  Almost everything is a text file.
\item
  The default app is not always the best.
\item
  Extensions are helpful, but not binding.
\item
  Useful to load and work with file in other folders.
\item
  Organize files logically into folders and subfolders.
\item
  Files and folders are the same whether accessed from the command line,
  or via apps.
\end{enumerate}
\end{frame}

\begin{frame}[fragile]{Commands covered (1)}
\protect\hypertarget{commands-covered-1}{}
\begin{Shaded}
\begin{Highlighting}[]
\BuiltInTok{pwd}
\BuiltInTok{cd}\NormalTok{ Downloads}
\BuiltInTok{cd}\NormalTok{ workshop}
\FunctionTok{ls} 
\FunctionTok{ls} \AttributeTok{{-}l}
\FunctionTok{ls} \AttributeTok{{-}l} \AttributeTok{{-}h}
\FunctionTok{ls} \AttributeTok{{-}lh} 
\FunctionTok{head}\NormalTok{ ted{-}sample.csv}
\FunctionTok{head}\NormalTok{ Command Line and Version Control.md}
\FunctionTok{head} \StringTok{"Command Line and Version Control.md"} 
\FunctionTok{head}\NormalTok{ Command}\DataTypeTok{\textbackslash{} }\NormalTok{Line}\DataTypeTok{\textbackslash{} }\NormalTok{and}\DataTypeTok{\textbackslash{} }\NormalTok{Version}\DataTypeTok{\textbackslash{} }\NormalTok{Control.md }
\FunctionTok{head} \AttributeTok{{-}n}\NormalTok{ 1 }\StringTok{"Command Line and Version Control.md"} 
\FunctionTok{head} \AttributeTok{{-}n}\NormalTok{ 10 }\StringTok{"Command Line and Version Control.md"} 
\FunctionTok{head} \AttributeTok{{-}n101}\NormalTok{ ted{-}sample.csv }\OperatorTok{\textgreater{}}\NormalTok{ ted100.csv}
\FunctionTok{head} \AttributeTok{{-}n101}\NormalTok{ ted{-}sample.csv }\OperatorTok{\textgreater{}\textgreater{}}\NormalTok{ ted100.csv}
\end{Highlighting}
\end{Shaded}
\end{frame}

\begin{frame}[fragile]{Commands covered (2)}
\protect\hypertarget{commands-covered-2}{}
\begin{Shaded}
\begin{Highlighting}[]
\FunctionTok{ls} \AttributeTok{{-}l} \PreprocessorTok{*}\NormalTok{.csv}
\FunctionTok{ls} \AttributeTok{{-}l} \PreprocessorTok{*}\NormalTok{sv}
\FunctionTok{mv}\NormalTok{ ted100.csv ted200.csv}
\FunctionTok{cp}\NormalTok{ ted200.csv ted200{-}backup.csv}
\FunctionTok{mv}\NormalTok{ ted200.csv ted200.py}
\FunctionTok{rm}\NormalTok{ ted200.py}
\FunctionTok{rm} \AttributeTok{{-}i}\NormalTok{ ted200{-}backup.csv }
\end{Highlighting}
\end{Shaded}
\end{frame}

\begin{frame}[fragile]{Version control}
\protect\hypertarget{version-control}{}
\begin{block}{We all know the problem}
\protect\hypertarget{we-all-know-the-problem}{}
\texttt{paper1.tex} \(\to\) \texttt{paper2.tex} \(\to\)
\texttt{paper\_final.tex} \(\to\) \texttt{paper\_final2.tex}
\end{block}

\begin{block}{We can have a system}
\protect\hypertarget{we-can-have-a-system}{}
\texttt{paper-2023-10-12.tex}
\end{block}

\begin{block}{Git is such a system}
\protect\hypertarget{git-is-such-a-system}{}
\texttt{git} is a tool to save snapshots of your files in a folder.
\end{block}
\end{frame}

\begin{frame}{Lessons from Version Control}
\protect\hypertarget{lessons-from-version-control}{}
\begin{enumerate}
\tightlist
\item
  Also useful on its own (use case: backup on a server without
  internet). But even better with GitHub, a collaborative website.
\item
  Commits are not automatic. You decide what to save and have to a give
  a commit message.
\item
  Merge conflicts are hard. But they are relatively rare.
\item
  Use cases for branches:

  \begin{itemize}
  \tightlist
  \item
    keep work separate to avoid intermittent conflicts
  \item
    review someone else's work
  \item
    experiment with big changes to code
  \end{itemize}
\end{enumerate}
\end{frame}

\begin{frame}[fragile]{Git commands covered}
\protect\hypertarget{git-commands-covered}{}
\begin{Shaded}
\begin{Highlighting}[]
\FunctionTok{git}\NormalTok{ init}
\FunctionTok{ls} \AttributeTok{{-}al}
\FunctionTok{git}\NormalTok{ status}
\FunctionTok{git}\NormalTok{ add ted200{-}backup.csv }
\FunctionTok{git}\NormalTok{ status}
\FunctionTok{git}\NormalTok{ config }\AttributeTok{{-}{-}global}\NormalTok{ user.name }\StringTok{"Miklos Koren"}
\FunctionTok{git}\NormalTok{ config }\AttributeTok{{-}{-}global}\NormalTok{ user.email }\StringTok{"foo@bar.com"}
\FunctionTok{git}\NormalTok{ commit }\AttributeTok{{-}m} \StringTok{"Add 200{-}line version of TED"}
\FunctionTok{git}\NormalTok{ log}
\FunctionTok{head} \AttributeTok{{-}n100}\NormalTok{ ted200{-}backup.csv }\OperatorTok{\textgreater{}}\NormalTok{ ted100{-}backup.csv}
\FunctionTok{mv}\NormalTok{ ted100{-}backup.csv ted200{-}backup.csv }
\FunctionTok{git}\NormalTok{ diff}
\FunctionTok{git}\NormalTok{ commit }\AttributeTok{{-}m} \StringTok{"Remove 102 lines from ted200{-}backup.csv"}
\FunctionTok{git}\NormalTok{ add ted200{-}backup.csv }
\FunctionTok{git}\NormalTok{ commit }\AttributeTok{{-}m} \StringTok{"Remove 102 lines from ted200{-}backup.csv"}
\FunctionTok{git}\NormalTok{ add }\StringTok{"Command Line and Version Control.md"} 
\FunctionTok{git}\NormalTok{ commit }\AttributeTok{{-}m} \StringTok{"Add workshop notes"}
\FunctionTok{git}\NormalTok{ diff}
\FunctionTok{git}\NormalTok{ commit }\AttributeTok{{-}am} \StringTok{"Add another lesson the command line part"}
\end{Highlighting}
\end{Shaded}
\end{frame}

\end{document}
